\section{Sectors}
\subsection{\tkzcname{tkzDrawSector}} 
\tkzHandBomb\  Attention the arguments vary according to the options.
\begin{NewMacroBox}{tkzDrawSector}{\oarg{local options}\parg{O,\dots}\parg{\dots}}%
\begin{tabular}{lll}%
options             & default & definition                         \\ 
\midrule
\TOline{towards}{towards}{$O$ is the center and the arc from $A$ to $(OB)$}
\TOline{rotate} {towards}{the arc starts from $A$ and the angle determines its length } 
\TOline{R}{towards}{We give the radius and two angles}
\TOline{R with nodes}{towards}{We give the radius and two points}
\bottomrule
\end{tabular} 

\medskip
You have to add, of course, all the styles of \TIKZ\ for tracings...

\medskip
\begin{tabular}{lll}%
\toprule
options             & arguments & example                         \\ 
\midrule
\TOline{towards}{\parg{pt,pt}\parg{pt}}{\tkzcname{tkzDrawSector(O,A)(B)}}
\TOline{rotate} {\parg{pt,pt}\parg{an}}{\tkzcname{tkzDrawSector[rotate,color=red](O,A)(90)}} 
\TOline{R}{\parg{pt,$r$}\parg{an,an}}{\tkzcname{tkzDrawSector[R,color=blue](O,2 cm)(30,90)}}
\TOline{R with nodes}{\parg{pt,$r$}\parg{pt,pt}}{\tkzcname{tkzDrawSector[R with nodes](O,2 cm)(A,B)}}
\bottomrule
\end{tabular}
\end{NewMacroBox}

Here are a few examples: 

\subsubsection{\tkzcname{tkzDrawSector} and \tkzname{towards}} 
There's no need to put \tkzname{towards}. You can use \tkzname{fill} as an option.

\begin{tkzexample}[latex=7cm,small]
\begin{tikzpicture}[scale=1]
  \tkzDefPoint(0,0){O}
  \tkzDefPoint(-30:3){A} 
  \tkzDefPointBy[rotation = center O angle -60](A) 
  \tkzDrawSector[fill=red!50](O,A)(tkzPointResult)
 \begin{scope}[shift={(-60:1cm)}]
  \tkzDefPoint(0,0){O}
  \tkzDefPoint(-30:3){A} 
  \tkzDefPointBy[rotation = center O angle -60](A) 
  \tkzDrawSector[fill=blue!50](O,tkzPointResult)(A)
  \end{scope}
\end{tikzpicture}   
\end{tkzexample}

\subsubsection{\tkzcname{tkzDrawSector} and \tkzname{rotate}}  
\begin{tkzexample}[latex=7cm,small]  
\begin{tikzpicture}[scale=2]
 \tkzDefPoint(0,0){O}
 \tkzDefPoint(2,2){A}
 \tkzDrawSector[rotate,draw=red!50!black,%
 fill=red!20](O,A)(30)
 \tkzDrawSector[rotate,draw=blue!50!black,%
 fill=blue!20](O,A)(-30)
\end{tikzpicture} 
\end{tkzexample}  

\subsubsection{\tkzcname{tkzDrawSector} and \tkzname{R}}  
\begin{tkzexample}[latex=7cm,small]
\begin{tikzpicture}[scale=1.25]
 \tkzDefPoint(0,0){O}
 \tkzDefPoint(2,-1){A}
 \tkzDrawSector[R,draw=white,%
 fill=red!50](O,2cm)(30,90)
 \tkzDrawSector[R,draw=white,%
 fill=red!60](O,2cm)(90,180)
 \tkzDrawSector[R,draw=white,%
 fill=red!70](O,2cm)(180,270)
 \tkzDrawSector[R,draw=white,%
 fill=red!90](O,2cm)(270,360) 
\end{tikzpicture}
\end{tkzexample}

\subsubsection{\tkzcname{tkzDrawSector} and \tkzname{R}}  
\begin{tkzexample}[latex=7cm,small]
\begin{tikzpicture}[scale=1.25]
 \tkzDefPoint(0,0){O}
 \tkzDefPoint(4,-2){A}
 \tkzDefPoint(4,1){B}
 \tkzDefPoint(3,3){C}
 \tkzDrawSector[R with nodes,%
                fill=blue!20](O,1 cm)(B,C)
 \tkzDrawSector[R with nodes,%
                fill=red!20](O,1.25 cm)(A,B)  
\tkzDrawSegments(O,A O,B O,C)
\tkzDrawPoints(O,A,B,C) 
\tkzLabelPoints(A,B,C) 
\tkzLabelPoints[left](O) 
\end{tikzpicture}
\end{tkzexample}

\subsubsection{\tkzcname{tkzDrawSector} and \tkzname{R with nodes}} 
\begin{tkzexample}[latex=7cm,small]
\begin{tikzpicture} [scale=.5]
 \tkzDefPoint(-1,-2){A}  
 \tkzDefPoint(1,3){B} 
 \tkzDefRegPolygon[side,sides=6](A,B) 
 \tkzGetPoint{O} 
 \tkzDrawPolygon[fill=black!10,
                 draw=blue](P1,P...,P6) 
 \tkzLabelRegPolygon[sep=1.05](O){A,...,F}
 \tkzDrawCircle[dashed](O,A)
 \tkzLabelSegment[above,sloped,
                  midway](A,B){\(A B = 16m\)}
 \foreach \i  [count=\xi from 1]  in {2,...,6,1}
   {%
    \tkzDefMidPoint(P\xi,P\i)
    \path (O) to [pos=1.1] node {\xi} (tkzPointResult) ;
    }
  \tkzDefRandPointOn[segment = P3--P5] 
  \tkzGetPoint{S}
  \tkzDrawSegments[thick,dashed,red](A,S S,B)
  \tkzDrawPoints(P1,P...,P6,S)
  \tkzLabelPoint[left,above](S){$S$}
  \tkzDrawSector[R with nodes,fill=red!20](S,2 cm)(A,B)
  \tkzLabelAngle[pos=1.5](A,S,B){$\alpha$}
\end{tikzpicture}
\end{tkzexample}

\subsection{\tkzcname{tkzFillSector}}
\tkzHandBomb\ Attention the arguments vary according to the options. 
\begin{NewMacroBox}{tkzFillSector}{\oarg{local options}\parg{O,\dots}\parg{\dots}}%
\begin{tabular}{lll}%
options          & default & definition      \\ 
\midrule
\TOline{towards}{towards}{$O$ is the center and the arc from $A$ to $(OB)$}
\TOline{rotate} {towards}{the arc starts from A and the angle determines its length } 
\TOline{R}{towards}{We give the radius and two angles}
\TOline{R with nodes}{towards}{We give the radius and two points}
\bottomrule
\end{tabular} 

\medskip
Of course, you have to add all the styles of \TIKZ\ for the tracings...

\medskip
\begin{tabular}{lll}%
\toprule
options             & arguments & example                         \\ 
\midrule
\TOline{towards}{\parg{pt,pt}\parg{pt}}{\tkzcname{tkzFillSector(O,A)(B)}}
\TOline{rotate} {\parg{pt,pt}\parg{an}}{\tkzcname{tkzFillSector[rotate,color=red](O,A)(90)}}
\TOline{R}{\parg{pt,$r$}\parg{an,an}}{\tkzcname{tkzFillSector[R,color=blue](O,2 cm)(30,90)}} 
\TOline{R with nodes}{\parg{pt,$r$}\parg{pt,pt}}{\tkzcname{tkzFillSector[R with nodes](O,2 cm)(A,B)}}
\end{tabular}   
\end{NewMacroBox} 

\subsubsection{\tkzcname{tkzFillSector} and \tkzname{towards}} 
It is useless to put \tkzname{towards} and you will notice that the contours are not drawn, only the surface is colored.
\begin{tkzexample}[latex=5.75cm,small]
\begin{tikzpicture}[scale=.6]
  \tkzDefPoint(0,0){O}
  \tkzDefPoint(-30:3){A} 
  \tkzDefPointBy[rotation = center O angle -60](A) 
  \tkzFillSector[fill=red!50](O,A)(tkzPointResult)
  \begin{scope}[shift={(-60:1cm)}]
   \tkzDefPoint(0,0){O}
   \tkzDefPoint(-30:3){A} 
   \tkzDefPointBy[rotation = center O angle -60](A) 
   \tkzFillSector[color=blue!50](O,tkzPointResult)(A)
  \end{scope}
\end{tikzpicture}
\end{tkzexample}


\subsubsection{\tkzcname{tkzFillSector} and \tkzname{rotate}} 
\begin{tkzexample}[latex=5.75cm,small]
\begin{tikzpicture}[scale=1.5]
 \tkzDefPoint(0,0){O} \tkzDefPoint(2,2){A}
 \tkzFillSector[rotate,color=red!20](O,A)(30)
 \tkzFillSector[rotate,color=blue!20](O,A)(-30)
\end{tikzpicture}    
\end{tkzexample} 

\subsection{\tkzcname{tkzClipSector}}
\tkzHandBomb\  Attention the arguments vary according to the options. 
\begin{NewMacroBox}{tkzClipSector}{\oarg{local options}\parg{O,\dots}\parg{\dots}}%
\begin{tabular}{lll}%
options             & default & definition                         \\ 
\midrule
\TOline{towards}{towards}{$O$ is the centre and the sector starts from $A$ to $(OB)$}
\TOline{rotate} {towards}{The sector starts from $A$ and the angle determines its amplitude. } 
\TOline{R}{towards}{We give the radius and two angles} 
\bottomrule
\end{tabular} 

\medskip
You have to add, of course, all the styles of \TIKZ\ for tracings...

\medskip   
\begin{tabular}{lll}%
\toprule
options             & arguments & example                         \\ 
\midrule
\TOline{towards}{\parg{pt,pt}\parg{pt}}{\tkzcname{tkzClipSector(O,A)(B)}}
\TOline{rotate} {\parg{pt,pt}\parg{angle}}{\tkzcname{tkzClipSector[rotate](O,A)(90)}} 
\TOline{R}{\parg{pt,$r$}\parg{angle 1,angle 2}}{\tkzcname{tkzClipSector[R](O,2 cm)(30,90)}}
\end{tabular}
\end{NewMacroBox}

\subsubsection{\tkzcname{tkzClipSector}} 
\begin{tkzexample}[latex=7cm,small] 
\begin{tikzpicture}[scale=1.5] 
  \tkzDefPoint(0,0){O}
  \tkzDefPoint(2,-1){A}
  \tkzDefPoint(1,1){B} 
  \tkzDrawSector[color=blue,dashed](O,A)(B)
  \tkzDrawSector[color=blue](O,B)(A)
  \tkzClipBB
  \begin{scope}
    \tkzClipSector(O,B)(A)
    \draw[fill=gray!20] (-1,0) rectangle (3,3);
  \end{scope}
  \tkzDrawPoints(A,B,O)
\end{tikzpicture} 
\end{tkzexample}

\endinput

